\section{Complete classifications for four, five and six parameters}\label{sec:classification}

\subsection{Method}
Let $G_m$ be the graph on $\Theta_m$ joining two cells at Hamming distance at most $\delta_m$. By Corollary~\ref{cor:pairs} every survivable set is a clique of $G_m$, so the maximal survivable sets are found inside the maximal cliques of $G_m$. The exact module (i) enumerates the maximal cliques (Bron--Kerbosch with pivoting \cite{BronKerbosch1973}), (ii) groups them into orbits of the hyperoctahedral group $B_m$ (coordinate permutations and sign flips, $|B_m|=2^mm!$) using the multiset of distance profiles as an invariant and confirming each class by an orbit-size count, (iii) solves the linear programme \eqref{eq:V} for one representative per orbit, (iv) inside each non-survivable representative searches for the maximal survivable subsets by branching only on the support of the certificate $q$ (every survivable subset omits a cell of $\supp q$), (v) tests the global maximality of each candidate against all one-cell extensions, (vi) certifies the thresholds of all subsets of each maximal representative (Section~\ref{sec:thresholds}), and (vii) maps the subsets through the group to obtain the complete families $\calS_m(L)$ as bit masks, from which the census, the maximal sets per level and the obstructions are read off. The pairwise-admissible family (all cliques of $G_m$) is obtained from the subsets of the maximal cliques. Runtimes are in Section~\ref{sec:float}.

\subsection{Statements}
\begin{theorem}[four parameters]\label{thm:m4}
$\calS_4(0)$ consists of the $16$ singletons; for every $L\ge1$, $\calS_4(L)$ consists of the $16$ singletons and the $32$ Hamming-adjacent pairs. No set of three or more cells is ever blind-survivable. The pairwise-admissible family has $48$ members, so from slack $1$ every pairwise-admissible set is survivable.
\end{theorem}
\begin{proof}
$\delta_4=1$ and $G_4$ is the four-cube graph, which is bipartite: its cliques are vertices and edges, so Corollary~\ref{cor:pairs} leaves singletons and adjacent pairs. Singletons have $\Lstar=0$ (constant matched action, $+3$); an adjacent pair has $V=3$ and first-step bound $1$ (every action costs one of the two cells at least $1$), and the alternating word from slack $1$ gives the trajectories $1,4,3,6,\dots$ and $1,0,3,2,\dots$, so $\Lstar=1$.
\end{proof}

\begin{theorem}[five parameters]\label{thm:m5}
$G_5$ has $192$ maximal cliques in three orbits: the $32$ radius-one balls (six cells), $80$ \emph{squares} (two-dimensional faces, distances $1,1,1,1,2,2$) and $80$ \emph{tetrahedra} (four cells pairwise at distance $2$). The balls have $\Lstar=0$, the squares $\Lstar=3$, the tetrahedra are never survivable; every clique of at most three cells lies in a ball. Hence $\calM_5(L)$ is the $32$ balls for $L=0,1,2$ and the $32$ balls together with the $80$ squares for $L\ge3$; $|\calS_5(L)|=1{,}616$ for $L\le2$ and $1{,}696$ for $L\ge3$, against $1{,}776$ pairwise-admissible sets, the difference being the $80$ tetrahedra.
\end{theorem}

\begin{theorem}[six parameters]\label{thm:m6}
$G_6$ has $10{,}752$ maximal cliques in seven orbits: $192$ of size $12$, $1{,}280$ of size $10$ and $9{,}280$ of size $8$ in five orbits ($160+2{,}880+3{,}840+1{,}920+480$). Exactly three orbits are survivable, and they are the maximal survivable sets:
\begin{enumerate}[leftmargin=2em,itemsep=1pt]
\item $192$ \emph{cylinders} $\Ball_5(c,1)\times\{\pm1\}$: the radius-one ball of a five-dimensional face, doubled along the sixth coordinate ($12$ cells, $V=3$, $\Lstar=1$); they attain Kleitman's bound $K(6)=12$.
\item $1{,}280$ \emph{boundary ten-sets}: in the three-dimensional face spanned by three chosen coordinates, the cube minus the vertex antipodal to a centre $c$ ($7$ cells), together with the three neighbours of $c$ along the remaining coordinates ($V=\tfr52$ exactly, $\Lstar=3$, first-step bound $1$).
\item $160$ \emph{three-dimensional subcubes} ($8$ cells, $V=3$, $\Lstar=5$).
\end{enumerate}
The remaining four orbits of maximal cliques have $V=2$ with the certificate $y=(0,0,0,0,1,-1)$ up to symmetry. The level census is: $|\calS_6(0)|=7{,}072$ with $\calM_6(0)$ the $64$ radius-one balls; $|\calS_6(1)|=|\calS_6(2)|=631{,}440$ with $\calM_6$ the $192$ cylinders and $1{,}280$ \emph{punctured subcubes} (a three-dimensional subcube minus one vertex, $\Lstar=1$); $|\calS_6(3)|=|\calS_6(4)|=774{,}800$ with $\calM_6$ the cylinders and the boundary ten-sets; $|\calS_6(L)|=774{,}960$ for $L\ge5$ with $\calM_6$ the cylinders, the ten-sets and the subcubes ($1{,}632$ maximal sets). The pairwise-admissible family has $853{,}680$ members; the $78{,}720$ admissible but non-survivable sets contain $2{,}400$ minimal obstructions, $1{,}440$ of four cells and $960$ of five, all with $V=2$. Among the survivable sets $143{,}360$ have $V=\tfr52$ exactly (all with $\Lstar=3$); the threshold histogram is $7{,}072$ sets at $0$, $624{,}368$ at $1$, $143{,}360$ at $3$ and $160$ at $5$.
\end{theorem}

The proofs of Theorems~\ref{thm:m5} and \ref{thm:m6} are the exact computation described above, with every linear programme certified on both sides, every threshold certified on both sides, and the orbit decomposition certified by orbit sizes; the records are in S3 and S4. Figure~\ref{fig:census} displays the census, and Table~\ref{tab:census} the counts. Three remarks on what the computation says.

\begin{table}[t]
\centering\small
\caption{Level census. For each $m$ and slack $L$ (initial stock $z_0=1+L/10$): the number of blind-survivable sets, the number and sizes of the maximal ones, the pairwise-admissible family (cliques of $G_m$) and the power set of the cells. The admissible family, not the power set, is the meaningful baseline: the pair law alone removes everything else.}
\label{tab:census}
\begin{tabular}{rrrlrr}
\toprule
$m$ & $L$ & $|\calS_m(L)|$ & maximal sets (size $\times$ count) & admissible & power set\\
\midrule
4 & 0 & 16 & $1\times16$ & 48 & $2^{16}=65{,}536$\\
4 & $\ge1$ & 48 & $2\times32$ & 48 & \\
\midrule
5 & 0--2 & 1{,}616 & $6\times32$ & 1{,}776 & $2^{32}\approx4.3\cdot10^{9}$\\
5 & $\ge3$ & 1{,}696 & $6\times32$, $4\times80$ & 1{,}776 & \\
\midrule
6 & 0 & 7{,}072 & $7\times64$ & 853{,}680 & $2^{64}\approx1.8\cdot10^{19}$\\
6 & 1--2 & 631{,}440 & $12\times192$, $7\times1{,}280$ & 853{,}680 & \\
6 & 3--4 & 774{,}800 & $12\times192$, $10\times1{,}280$ & 853{,}680 & \\
6 & $\ge5$ & 774{,}960 & $12\times192$, $10\times1{,}280$, $8\times160$ & 853{,}680 & \\
\bottomrule
\end{tabular}
\end{table}

\begin{figure}[t]
\centering
\includegraphics[width=\textwidth]{fig_census.pdf}
\caption{The census of the six-cube by level: survivable sets (left axis, thousands) and maximal sets by type (right axis). The jumps at $L=1$, $3$ and $5$ are the thresholds of the cylinders and punctured subcubes, the boundary ten-sets, and the full three-dimensional subcubes. The dashed line is the pairwise-admissible family.}
\label{fig:census}
\end{figure}

\paragraph{Kleitman's bound is attained but not by balls.}
At $m=4$ and $m=6$ ($\delta_m$ odd) the extremal sets of Kleitman's theorem are doubled balls of a facet; at $m=5$ ($\delta_m$ even) they are balls. In all three cases the extremal sets are survivable, and at $m=4,6$ from slack $1$ rather than at the floor, because the matched action of one centre costs the cells of the other layer one tenth; Corollary~\ref{cor:kleitman} shows that this holds for every $m$.

\paragraph{The largest sets are not the hardest.}
The thresholds are not monotone in size: the twelve-cell cylinders need slack $1$, the ten-cell boundary sets slack $3$, the eight-cell subcubes slack $5$, and the seven-cell punctured subcubes slack $1$ again. The subcube is expensive because its $2^3$ cells disagree on three coordinates, so every sign action is at distance $3$ from one of them ($D=-5$), whereas the cylinder's cells disagree on one coordinate each with a common five-dimensional centre.

\paragraph{Uniform horizon.}
Every minimal obstruction of the three cubes---far pairs at distance $\delta_m+1$, the tetrahedra of the five-cube, the $2{,}400$ four- and five-cell obstructions of the six-cube---has $V=2$, so Corollary~\ref{cor:horizon} gives the same horizon for all of them.

\begin{corollary}[uniform horizon for $m\le6$]\label{cor:uniform}
For $m\in\{4,5,6\}$, every set that is not blind-survivable loses a cell within $L=10(z_0-1)$ steps from common slack $L$, whatever the blind word; the bound is attained by far pairs at $m=4,6$.
\end{corollary}
\begin{proof}
A non-survivable set contains a minimal obstruction $O$ (Theorem~\ref{thm:criterion}, support bound), and $V(O)=2$ for every minimal obstruction by the enumeration, so $5-2V(O)=1$ and the horizon is $\lfloor L/1\rfloor=L$.
\end{proof}

The corollary gives an exact meaning to the ``bands'' of the earlier editions and to the finite-horizon evaluations of Section~\ref{sec:pbvi}: a set that survives $L+1$ steps from slack $L$ is survivable for ever.
